% Fases 2 y 3 · Modelado y evaluación, y objetivos secundarios
\begin{tikzpicture}
% Estilos compartidos de los diagramas de la arquitectura (figuras de la sección Visión general)
\tikzset{
  font=\sffamily\scriptsize,
  comp/.style={draw=uoc-azul, line width=0.6pt, fill=papel, rounded corners=2pt,
               align=center, minimum height=10mm, text width=19mm, inner sep=2pt},
  dato/.style={draw=serie-2, line width=0.6pt, fill=uoc-cian-suave, shape=cylinder,
               shape border rotate=90, aspect=0.15, align=center,
               minimum height=11mm, minimum width=19mm, inner sep=1pt},
  sec/.style={dashed, draw=serie-base, text=tinta-suave},
  flecha/.style={-{Latex[length=2mm]}, line width=0.6pt, draw=uoc-azul},
  flechasec/.style={flecha, draw=serie-base, dashed},
  carril/.style={draw=none, fill=superficie, rounded corners=3pt, inner sep=3mm},
  rotulo/.style={font=\sffamily\scriptsize\bfseries, text=uoc-azul, anchor=north west, inner sep=2pt},
  num/.style={circle, fill=uoc-azul, text=papel, inner sep=0pt, minimum size=4.2mm, font=\sffamily\scriptsize\bfseries}
}

% --- Conjunto unificado (viene de la figura anterior) ---
\node[dato] (unificado) at (1.4,1.5) {Conjunto\\unificado};

% --- Objetivo secundario: variables cartográficas ---
\node[dato, sec, fill=papel, minimum width=27mm] (carto) at (1.9,3.2) {Cartografía pública\\[-1pt]{\tiny PNOA · Catastro · MDT02}};
\node[comp, sec] (obtencion)   at (5.4,3.2) {OS2 · Variables cartográficas};
\node[comp, sec] (unificacion) at (5.4,1.5) {Unificación};
\node[dato, sec, fill=papel, minimum width=30mm] (final) at (9.2,1.5) {Datos unificados con\\variables cartográficas};
\draw[flechasec] (carto) -- (obtencion);
\draw[flechasec] (obtencion) -- (unificacion);
\draw[flechasec] (unificado) -- (unificacion);
\draw[flechasec] (unificacion) -- (final);

% --- Modelado y evaluación ---
\node[comp] (variables)  at (1.4,-0.8) {Conjuntos de variables según origen};
\node[comp] (division)   at (4.0,-0.8) {División entrenamiento / test};
\node[comp] (entrena)    at (6.6,-0.8) {Entrenamiento de modelos};
\node[comp] (evalua)     at (9.2,-0.8) {Evaluación};
\node[comp] (interpreta) at (11.8,-0.8) {Interpretación y contraste};
\draw[flecha] ([xshift=-3mm]unificado.south) -- ([xshift=-3mm]variables.north);
\draw[flechasec] (final.south) -- ++(0,-0.6) -| ([xshift=3mm]variables.north);
\draw[flecha] (variables) -- (division);
\draw[flecha] (division) -- (entrena);
\draw[flecha] (entrena) -- (evalua);
\draw[flecha] (evalua) -- (interpreta);

\node[draw=serie-2, line width=0.6pt, fill=uoc-cian-suave, shape=tape, tape bend top=none, tape bend height=1.5mm, align=center, minimum height=11mm, minimum width=19mm, inner sep=2pt] (resultados) at (9.2,-2.6) {Resultados\\e informes};
\node[comp, sec] (os1) at (5.8,-2.6) {OS1 · Visor GIS};
\draw[flecha] (evalua) -- (resultados);
\draw[flecha] (interpreta.south) |- (resultados.east);
\draw[flechasec] (resultados) -- (os1);

\foreach \n/\c in {2/variables, 3/division, 4/entrena, 5/evalua, 6/interpreta}
  \node[num] at (\c.north west) {\n};

\begin{scope}[on background layer]
  \node[carril, fit=(variables)(interpreta)(resultados)(os1)] (cmodelo) {};
\end{scope}
\node[rotulo, anchor=south west] at (cmodelo.south west) {Modelado y evaluación};
\end{tikzpicture}
